\documentclass[12pt]{amsart}
\usepackage[margin=1in]{geometry}
\usepackage[utf8]{inputenc}
\usepackage{amsfonts}
\usepackage{amsthm}
\usepackage{amsmath}
\usepackage{amssymb}
\usepackage{graphicx}
\usepackage{enumitem}
\usepackage{quiver}
\usepackage{mathrsfs}
\usepackage{mathtools}
\usepackage{hyperref}
\usepackage{cleveref}
\hypersetup{
    colorlinks=true,
    linkcolor=blue,
    filecolor=green,
    urlcolor=cyan,
    citecolor=magenta,
    pdftitle={Hodge Loci, Mumford--Tate Groups, and the Hodge Conjecture for Abelian Varieties},
    pdfpagemode=FullScreen,
}
\urlstyle{same}


\newcommand{\R}{\mathbb{R}}
\newcommand{\Q}{\mathbb{Q}}
\newcommand{\C}{\mathbb{C}}
\newcommand{\Z}{\mathbb{Z}}
\newcommand{\PP}{\mathbb{P}}
\newcommand{\Gm}{\mathbb{G}_m}
\newcommand{\Del}{\mathbb{S}}
\newcommand{\End}{\text{End}}
\newcommand{\Aut}{\text{Aut}}
\newcommand{\Hom}{\text{Hom}}
\newcommand{\id}{\text{id}}

\DeclareMathOperator{\GL}{GL}
\DeclareMathOperator{\SL}{SL}
\DeclareMathOperator{\Sp}{Sp}
\DeclareMathOperator{\GSp}{GSp}
\DeclareMathOperator{\MT}{MT}
\DeclareMathOperator{\Hg}{Hg}
\DeclareMathOperator{\Hdg}{Hdg}
\DeclareMathOperator{\CH}{CH}
\DeclareMathOperator{\NL}{NL}
\DeclareMathOperator{\HL}{HL}
\DeclareMathOperator{\Res}{Res}
\DeclareMathOperator{\Stab}{Stab}
\DeclareMathOperator{\Nm}{Nm}
\DeclareMathOperator{\Sym}{Sym}
\DeclareMathOperator{\Pic}{Pic}
\DeclareMathOperator{\Tr}{Tr}
\DeclareMathOperator{\divisor}{div}

\newcommand{\cA}{\mathcal{A}}
\newcommand{\cO}{\mathcal{O}}
\newcommand{\fH}{\mathfrak{H}}
\newcommand{\fh}{\mathfrak{h}}
\newcommand{\fsl}{\mathfrak{sl}}
\newcommand{\fsp}{\mathfrak{sp}}
\newcommand{\ol}{\overline}
\newcommand{\wt}{\widetilde}

\newtheorem{theorem}{Theorem}[section]
\newtheorem{proposition}[theorem]{Proposition}
\newtheorem{lemma}[theorem]{Lemma}
\newtheorem{corollary}[theorem]{Corollary}
\newtheorem{conjecture}[theorem]{Conjecture}
\newtheorem{fact}[theorem]{Fact}

\theoremstyle{definition}
\newtheorem{definition}[theorem]{Definition}
\newtheorem{example}[theorem]{Example}

\theoremstyle{remark}
\newtheorem{remark}[theorem]{Remark}

\numberwithin{equation}{section}


\title[Hodge Loci, MT Groups, and the Hodge Conjecture]{Hodge Loci, Mumford--Tate Groups, and the Hodge Conjecture for Abelian Varieties}
\author{Edward Varvak}

\begin{document}

\maketitle

These are notes for an introductory talk on Hodge (or Noether--Lefschetz) loci,
Mumford--Tate groups, and the Hodge conjecture, with abelian varieties as the
running example. The guiding principle is
\begin{center}
    \emph{Hodge classes on $A$ $=$ invariants of the Mumford--Tate group of $A$,}
\end{center}
so that computing Mumford--Tate groups computes Hodge classes. Hodge loci are
where Hodge classes ``jump'', and on the moduli space of abelian varieties they
correspond to algebraic subgroups. We end by applying this to the fibers of a
maximally varying Lagrangian fibration (\Cref{thm:main,cor:main}).

Throughout, varieties are over $\C$, and cohomology means singular cohomology of
the associated complex analytic space.

\section{Algebraic cycles and the Hodge conjecture}

Let $X$ be a smooth projective variety of dimension $n$.

\begin{definition}
    An \emph{algebraic cycle of codimension $p$} on $X$ is a finite formal sum
    $Z = \sum_i n_i Z_i$ with $n_i \in \Z$ and $Z_i \subseteq X$ irreducible
    closed subvarieties of codimension $p$. These form a free abelian group
    $\mathcal{Z}^p(X)$.
\end{definition}

\begin{definition}
    A cycle $Z \in \mathcal{Z}^p(X)$ is \emph{rationally equivalent to zero} if
    there are subvarieties $W_j \subseteq X$ of codimension $p-1$ and rational
    functions $f_j \in \C(W_j)^\times$ with $Z = \sum_j \divisor(f_j)$. The
    \emph{Chow group} is $\CH^p(X) = \mathcal{Z}^p(X)/\sim_{\mathrm{rat}}$.
\end{definition}

For example $\CH^1(X) = \Pic(X)$. Equivalently, $Z \sim_{\mathrm{rat}} 0$ if
$Z$ is a sum of differences $\Gamma_0 - \Gamma_\infty$ of fibers of
subvarieties $\Gamma \subseteq X \times \PP^1$ flat over $\PP^1$, so rationally
equivalent cycles are deformations of each other over $\PP^1$.

\subsection{The cycle class map}
If $Z \subseteq X$ is irreducible of codimension $p$, integration over the
smooth locus of $Z$ gives a well-defined functional
$\eta \mapsto \int_Z \eta$ on $H^{2n-2p}(X, \C)$ (Lelong), and hence by
Poincar\'e duality a class $[Z] \in H^{2p}(X, \C)$. In fact
$[Z] \in H^{2p}(X, \Z)$: it is the Poincar\'e dual of the fundamental class of
$Z$. Since cohomology classes are locally constant in families, this descends
to the \emph{cycle class map}
\begin{equation*}
    \mathrm{cl} \colon \CH^p(X) \to H^{2p}(X, \Z).
\end{equation*}

\begin{lemma}\label{lem:cycle-class-hodge}
    For any cycle $Z$, the image of $[Z]$ in $H^{2p}(X, \C)$ lies in
    $H^{p,p}(X)$.
\end{lemma}
\begin{proof}
    By Poincar\'e duality it suffices to show $\int_Z \eta = 0$ for every closed
    form $\eta$ of type $(a, b)$ with $a + b = 2n - 2p$ and
    $(a, b) \neq (n - p, n - p)$. But then $a > n - p$ or $b > n - p$, and an
    $(a,b)$-form restricts to zero on the complex manifold $Z^{\mathrm{sm}}$,
    which has dimension $n - p$.
\end{proof}

\begin{definition}
    The groups of \emph{integral} and \emph{rational Hodge classes} are
    \begin{align*}
        \Hdg^p(X, \Z) &= \{\alpha \in H^{2p}(X, \Z) : \alpha_\C \in H^{p,p}(X)\},\\
        \Hdg^p(X) = \Hdg^p(X, \Q) &= H^{2p}(X, \Q) \cap H^{p,p}(X).
    \end{align*}
    A Hodge class is \emph{algebraic} if it lies in the image of $\mathrm{cl}$
    (respectively of $\mathrm{cl} \otimes \Q$).
\end{definition}

\begin{remark}
    The Hodge conjecture only concerns the \emph{image} of $\mathrm{cl}$, so
    rational equivalence plays no real role below. It is just the most natural
    equivalence relation to put on cycles.
\end{remark}

\subsection{The (integral) Hodge conjecture}

\begin{conjecture}[Integral Hodge conjecture]
    For all $p$, the cycle class map
    \begin{equation*}
        \mathrm{cl} \colon \CH^p(X) \to \Hdg^p(X, \Z)
    \end{equation*}
    is surjective.
\end{conjecture}

\begin{theorem}[Lefschetz $(1,1)$]
    The integral Hodge conjecture holds for $p = 1$: every class in
    $\Hdg^1(X, \Z)$ is the first Chern class of a line bundle.
\end{theorem}

It also holds trivially for $p = 0$ and $p = n$. Unfortunately, it is false in
general.

\begin{theorem}[Atiyah--Hirzebruch \cite{AH62}]
    There are smooth projective $X$ with torsion classes in $H^4(X, \Z)$ that
    are not algebraic. (Torsion classes vanish in $H^4(X, \C)$, so they are
    automatically Hodge classes.)
\end{theorem}

\begin{theorem}[Koll\'ar \cite{Kollar92}]
    Let $X \subseteq \PP^4$ be a very general hypersurface of degree $d$, where
    $p^3 \mid d$ for some prime $p \geq 5$. Then every curve on $X$ has degree
    divisible by $p$.
\end{theorem}

By the Lefschetz hyperplane theorem and Poincar\'e duality,
$H^4(X, \Z) \cong \Z$ is generated by a class $\ell$ with $\ell \cdot h = 1$,
and every class in $H^4(X,\Z)$ is Hodge. So $\ell$ is a non-torsion Hodge class
that is not algebraic, even though $d \ell = h^2$ is. The standard fix is to
pass to $\Q$-coefficients.

\begin{conjecture}[Hodge conjecture]
    For all $p$, $\Hdg^p(X)$ is spanned over $\Q$ by classes of algebraic
    cycles, i.e.\ $\mathrm{cl} \otimes \Q \colon \CH^p(X)_\Q \to \Hdg^p(X)$ is
    surjective.
\end{conjecture}

\begin{proposition}
    The Hodge conjecture holds in degrees $p = 0, 1, n-1, n$. In particular it
    holds for all $X$ of dimension at most $3$.
\end{proposition}
\begin{proof}
    The cases $p = 0, n$ are trivial and $p = 1$ is Lefschetz $(1,1)$. For
    $p = n - 1$, let $L$ be cup product with the class of a hyperplane section
    $H$. By hard Lefschetz,
    $L^{n-2} \colon H^2(X, \Q) \to H^{2n-2}(X, \Q)$ is an isomorphism, and it
    maps $H^{1,1}$ onto $H^{n-1,n-1}$, so it gives an isomorphism
    $\Hdg^1(X) \cong \Hdg^{n-1}(X)$. It sends $[D]$ to
    $[D \cap H_1 \cap \cdots \cap H_{n-2}]$ for general hyperplane sections
    $H_i$, so it preserves algebraic classes.
\end{proof}

\begin{remark}\label{rmk:AV-intro}
    If $A$ is an abelian variety, then $H^*(A, \Z)$ is an exterior algebra on
    $H^1(A, \Z)$, and in particular torsion-free. So there are no
    Atiyah--Hirzebruch type counterexamples on abelian varieties. Nevertheless,
    the integral Hodge conjecture fails for very general principally polarized
    abelian varieties of dimension $\geq 4$ by recent work of Engel, de Gaay
    Fortman, and Schreieder; see \Cref{thm:EdGFS} below.
\end{remark}

\section{Abelian varieties and their moduli}

\subsection{Hodge structures}
Let $\Del = \Res_{\C/\R} \Gm$ be the Deligne torus, so $\Del(\R) = \C^\times$.
A \emph{$\Q$-Hodge structure of weight $w$} is a finite-dimensional
$\Q$-vector space $V$ together with a homomorphism of real algebraic groups
$h \colon \Del \to \GL(V_\R)$ such that
\begin{equation*}
    V_\C = \bigoplus_{p + q = w} V^{p,q}, \qquad
    h(z) \text{ acts on } V^{p,q} \text{ by } z^p \ol{z}^q.
\end{equation*}
(Conventions differ by a sign, which will not matter.) Two consequences we will
use constantly:
\begin{itemize}
    \item for $r \in \R^\times$, $h(r) = r^w$ is a scalar;
    \item for $z \in U(1) = \{|z| = 1\}$, $h(z)$ acts on $V^{p,q}$ by
    $z^{p-q}$.
\end{itemize}
Duals, tensor products, and exterior powers of Hodge structures are Hodge
structures; we write
$T^{a,b}(V) = V^{\otimes a} \otimes (V^\vee)^{\otimes b}$, of weight
$(a-b)w$, and call $T^{a,b}(V)$, $\Lambda^k V$, $\Sym^k V$ and sums of these
\emph{tensor constructions} of $V$. If $W$ has weight $2p$, a \emph{Hodge
class} in $W$ is an element of $W \cap W^{p,p}$.

\subsection{Abelian varieties}
A complex abelian variety of dimension $g$ is a projective complex torus
$A = \C^g / \Lambda$. We always write
\begin{equation*}
    V = H^1(A, \Q).
\end{equation*}
Then $H^k(A, \Z) = \Lambda^k H^1(A, \Z)$,
$H^1(A, \C) = H^{1,0} \oplus H^{0,1}$ with $H^{1,0} = H^0(A, \Omega^1_A)$, and
$H^{p,q}(A) = \Lambda^p H^{1,0} \otimes \Lambda^q H^{0,1}$. So:

\begin{center}
    \emph{The Hodge structure on $H^*(A, \Q) = \Lambda^\bullet V$ is determined
    by the weight one Hodge structure $V$,}
\end{center}
and $\Hdg^p(A)$ is the set of Hodge classes in $\Lambda^{2p} V$.

A polarization is the class $\theta = c_1(L) \in H^2(A, \Z)$ of an ample line
bundle. We will only consider principal polarizations, for which there is a
basis $e_1, \ldots, e_g, f_1, \ldots, f_g$ of $H^1(A, \Z)$ with
\begin{equation*}
    \theta = \sum_{i=1}^g e_i \wedge f_i \in \Lambda^2 H^1(A, \Z).
\end{equation*}
Let $\psi$ be the symplectic form on $V$ for which $(e_i, f_i)$ is a symplectic
basis. We write
\begin{equation*}
    \Sp(V) = \Sp(V, \psi) = \Stab_{\GL(V)}(\theta) \cong \Sp_{2g, \Q}, \qquad
    \GSp(V) = \Gm \cdot \Sp(V),
\end{equation*}
the latter being the group of $g \in \GL(V)$ that preserve $\psi$ up to a
scalar.

\subsection{The Siegel upper half space and $\cA_g$}
Fix $V_\Z = \Z^{2g}$ with its standard symplectic form $\psi$, and let
$V = V_\Z \otimes \Q$. Let
\begin{equation*}
    \fH_g = \left\{ h \colon \Del \to \GL(V_\R) \;\middle|\;
    \begin{array}{l}
        h \text{ is a Hodge structure of weight } 1 \text{ with}\\
        \psi(h(i) x, h(i) y) = \psi(x, y), \text{ and } \psi(x, h(i)y)
        \text{ positive definite}
    \end{array}\right\}.
\end{equation*}
The last two conditions say that $\psi$ is a polarization of the Hodge
structure (for a suitable sign convention).

\begin{theorem}[Riemann]
    Points of $\fH_g$ are in bijection with isomorphism classes of pairs
    $(A, \beta)$, where $A$ is a principally polarized abelian variety of
    dimension $g$ and $\beta \colon H^1(A, \Z) \cong V_\Z$ is a symplectic
    isomorphism.
\end{theorem}

The group $\Sp(V_\R) = \Sp_{2g}(\R)$ acts on $\fH_g$ by conjugation,
$g \cdot h = g h g^{-1}$, and this action is transitive with stabilizer
isomorphic to $U(g)$. So
\begin{equation*}
    \fH_g \cong \Sp_{2g}(\R)/U(g) \cong
    \{\tau \in M_g(\C) : \tau = \tau^T, \ \mathrm{Im}(\tau) > 0\},
\end{equation*}
the classical Siegel upper half space, of complex dimension $g(g+1)/2$.
Forgetting the marking $\beta$, we get the moduli space of principally polarized
abelian varieties
\begin{equation*}
    \cA_g = \Sp_{2g}(\Z) \backslash \fH_g,
\end{equation*}
which is a quasi-projective variety of dimension $g(g+1)/2$ (Baily--Borel).

\begin{remark}[Shimura varieties]
    $\cA_g$ is the basic example of a \emph{Shimura variety}: a quotient
    $\Gamma \backslash X$, where $G$ is a reductive group over $\Q$, $X$ is a
    $G(\R)$-conjugacy class of homomorphisms $\Del \to G_\R$ (satisfying axioms
    that make $X$ a Hermitian symmetric domain), and $\Gamma \subseteq G(\Q)$
    is an arithmetic subgroup. Here $G = \GSp(V)$ and $X = \fH_g$. The point
    is that the period domain is a single orbit of a real algebraic group
    acting by conjugation on Hodge structures. We will see that the Hodge loci
    in $\cA_g$ have the same form, for smaller groups (\Cref{sec:hodge-loci}).
\end{remark}

The key feature of $\fH_g$ is that for every $h \in \fH_g$,
\begin{equation*}
    H^{2p}(A_h, \Q) = \Lambda^{2p} V
\end{equation*}
is the \emph{same} $\Q$-vector space; only the Hodge decomposition, induced by
$h$, moves. This makes Hodge loci in $\fH_g$ very concrete.

\section{Mumford--Tate groups}

\begin{definition}
    Let $V$ be a $\Q$-Hodge structure given by $h \colon \Del \to \GL(V_\R)$.
    \begin{itemize}
        \item The \emph{Mumford--Tate group} $\MT(V)$ is the smallest
        $\Q$-algebraic subgroup $G \subseteq \GL(V)$ with
        $h(\C^\times) \subseteq G(\R)$.
        \item The \emph{Hodge group} (or special Mumford--Tate group) $\Hg(V)$
        is the smallest $\Q$-algebraic subgroup $G \subseteq \GL(V)$ with
        $h(U(1)) \subseteq G(\R)$.
    \end{itemize}
    For an abelian variety we write $\MT(A) = \MT(H^1(A, \Q))$ and
    $\Hg(A) = \Hg(H^1(A,\Q))$.
\end{definition}

\begin{remark}\label{rmk:MT-basic}
    \begin{enumerate}[label=(\arabic*)]
        \item $\MT(V)$ and $\Hg(V)$ are connected, since the identity component
        also satisfies the defining property.
        \item If $V$ has weight $w \neq 0$, then $\MT(V) = \Gm \cdot \Hg(V)$,
        since $\C^\times = \R_{>0} \cdot U(1)$ and $h(r) = r^w$ is a scalar.
        \item If $V$ is polarizable (e.g.\ $V = H^1(A,\Q)$), then $\MT(V)$ and
        $\Hg(V)$ are reductive \cite{Deligne82, Moonen04}.
        \item For an abelian variety, $\Hg(A^m) = \Hg(A)$, acting diagonally on
        $H^1(A^m, \Q) = V^{\oplus m}$: the diagonal copy of $\Hg(A)$ contains
        $h(U(1))$, and any $\Q$-subgroup of it containing $h(U(1))$ projects
        onto $\Hg(A)$.
    \end{enumerate}
\end{remark}

The reason to care about $\Hg$ is the following.

\begin{proposition}\label{prop:key}
    Let $W$ be a tensor construction of $V$ of even weight $2p$. Then
    $t \in W$ is a Hodge class if and only if $t$ is fixed by $\Hg(V)$. In
    particular, for an abelian variety $A$,
    \begin{equation*}
        \Hdg^p(A) = (\Lambda^{2p} V)^{\Hg(A)}.
    \end{equation*}
\end{proposition}
\begin{proof}
    Since $\Hg(V) \subseteq \GL(V)$, it acts on every tensor construction of
    $V$. An element $z \in U(1)$ acts on $W^{a,b}$ by $z^{a-b}$, so
    $t \in W^{p,p}$ if and only if $h(z) t = t$ for all $z \in U(1)$. Since $t$
    is rational, its stabilizer $\Stab(t) \subseteq \GL(V)$ is a $\Q$-algebraic
    subgroup, and so
    \begin{equation*}
        t \in W^{p,p} \iff h(U(1)) \subseteq \Stab(t)(\R) \iff
        \Hg(V) \subseteq \Stab(t)
    \end{equation*}
    by minimality of $\Hg(V)$.
\end{proof}

\begin{remark}
    The Mumford--Tate group itself acts on a Hodge class of weight $2p$ through
    the character $\Gm \to \Gm$, $r \mapsto r^{2p}$ on the central $\Gm$, which
    is why it is more convenient to work with $\Hg$. (Alternatively, Tate twist
    everything to weight $0$.)
\end{remark}

\begin{corollary}\label{cor:key}
    Let $A$ be a principally polarized abelian variety, $V = H^1(A, \Q)$.
    \begin{enumerate}[label=(\alph*)]
        \item Since $\theta \in \Lambda^2 V$ is a Hodge class,
        $\Hg(A) \subseteq \Sp(V)$ and $\MT(A) \subseteq \GSp(V)$.
        \item Sub-Hodge structures of $V$ are exactly the $\Hg(A)$-stable
        subspaces.
        \item Morphisms of Hodge structures are Hodge classes in
        $\Hom(V, V') = V^\vee \otimes V'$ (weight $0$). So
        $\End_{\mathrm{HS}}(V) = \End(V)^{\Hg(A)}$, and more generally, since
        $\Hom(A, B) \otimes \Q \cong \Hom_{\mathrm{HS}}(H^1(B,\Q), H^1(A,\Q))$,
        endomorphisms and isogenies of abelian varieties are detected by $\Hg$.
    \end{enumerate}
\end{corollary}

In fact $\Hg(V)$ is \emph{determined} by the Hodge classes in all tensor
constructions:

\begin{fact}[Deligne \cite{Deligne82}]\label{fact:deligne}
    If $V$ is polarizable, then $\Hg(V)$ is the subgroup of $\GL(V)$ fixing
    every Hodge class in every $T^{a,b}(V)$.
\end{fact}

\subsection{Elliptic curves}
For an elliptic curve $E$, $V = H^1(E, \Q)$ is $2$-dimensional, so
$\Sp(V) = \SL(V)$.

\begin{proposition}\label{prop:elliptic}
    Let $E$ be an elliptic curve and $V = H^1(E, \Q)$.
    \begin{enumerate}[label=(\roman*)]
        \item If $E$ does not have complex multiplication, then
        $\Hg(E) = \SL(V) \cong \SL_2$ and $\MT(E) = \GL(V) \cong \GL_2$.
        \item If $\End(E) \otimes \Q = K$ is an imaginary quadratic field, then
        $V$ is a $1$-dimensional $K$-vector space,
        \begin{equation*}
            \MT(E) = T_K := \Res_{K/\Q} \Gm, \qquad
            \Hg(E) = U_K := \ker\left(\Nm \colon T_K \to \Gm\right).
        \end{equation*}
    \end{enumerate}
\end{proposition}
\begin{proof}
    $\Hg = \Hg(E)$ is a connected reductive subgroup of $\SL(V)$, and it is
    nontrivial since $h(z)$ acts on $V^{1,0}$ by $z$. The connected reductive
    subgroups of $\SL_2$ are $1$, the $1$-dimensional tori, and $\SL_2$
    itself.

    Suppose $\Hg = T$ is a torus. Over $\C$, $T$ acts on $V_\C$ through two
    distinct characters $\chi, \chi^{-1}$ (the only scalars in $\SL_2$ are
    $\pm 1$), so its centralizer $K' = \End(V)^T$ is a $2$-dimensional
    commutative $\Q$-algebra. By \Cref{cor:key}(c),
    $K' \subseteq \End(E) \otimes \Q$. Since $\End(E) \otimes \Q$ is either
    $\Q$ or an imaginary quadratic field, $E$ has CM.

    Conversely, suppose $E$ has CM by $K$. Then $K$ acts on $V$ commuting with
    $\Hg$, so $\Hg$ lies in the centralizer of $K$ in $\SL(V)$, which is
    $T_K \cap \SL(V) = U_K$. This is a $1$-dimensional torus, and $\Hg$ is
    nontrivial and connected, so $\Hg = U_K$. Finally
    $\MT(E) = \Gm \cdot U_K = T_K$.
\end{proof}

\begin{remark}
    $E_\tau = \C/(\Z + \tau\Z)$ has CM if and only if $\tau \in \fH_1$ is an
    imaginary quadratic number. So CM elliptic curves correspond to a
    \emph{countable dense} subset of $\fH_1$.
\end{remark}

\begin{example}[Hodge classes on $E \times E$]\label{ex:ExE}
    We have $H^1(E \times E) = V \oplus V$, and by \Cref{rmk:MT-basic}(4),
    $\Hg(E \times E) = \Hg(E)$ acting diagonally. Then
    \begin{equation*}
        H^2(E \times E, \Q) = \Lambda^2 V \oplus (V \otimes V) \oplus
        \Lambda^2 V.
    \end{equation*}
    $\Hg(E) \subseteq \SL(V)$ acts trivially on each $\Lambda^2 V$, giving the
    classes of $E \times 0$ and $0 \times E$. Using the
    $\GL(V)$-equivariant isomorphism $V \cong V^\vee \otimes \Lambda^2 V$, we
    get $V \otimes V \cong \End(V)$ as $\Hg(E)$-representations, so
    $(V \otimes V)^{\Hg(E)} \cong \End(V)^{\Hg(E)} = \End(E) \otimes \Q$.
    Therefore
    \begin{equation*}
        \rho(E \times E) = 2 + \dim_\Q \End(E) \otimes \Q =
        \begin{cases}
            3 & E \text{ without CM,}\\
            4 & E \text{ with CM.}
        \end{cases}
    \end{equation*}
    The extra classes are the diagonal and, in the CM case, the graph of a
    complex multiplication. So ``$E$ has CM'' is detected by a Hodge class on
    $E \times E$.
\end{example}

\begin{proposition}\label{prop:product-elliptic}
    Let $E_1, \ldots, E_n$ be pairwise non-isogenous elliptic curves without
    CM, $A = E_1 \times \cdots \times E_n$, and $V_i = H^1(E_i, \Q)$. Then
    \begin{align*}
        \Hg(A) &= \SL(V_1) \times \cdots \times \SL(V_n),\\
        \MT(A) &= \Gm \cdot \SL_2^n = \Big\{(g_i) \in \prod_i \GL(V_i) :
        \det g_1 = \cdots = \det g_n\Big\}.
    \end{align*}
\end{proposition}
\begin{proof}
    $\prod_i \Hg(E_i) = \prod_i \SL(V_i)$ is a $\Q$-group containing
    $h(U(1)) = \{(h_1(z), \ldots, h_n(z))\}$, so $\Hg(A)$ is contained in it.
    The projection of $\Hg(A)$ to $\SL(V_i)$ is a $\Q$-group containing
    $h_i(U(1))$, so it is all of $\SL(V_i)$ by \Cref{prop:elliptic}. Let
    $\fh \subseteq \bigoplus_i \fsl(V_i)$ be the Lie algebra of $\Hg(A)$.

    Take $n = 2$. Then $\fh \cap (\fsl(V_1) \oplus 0)$ is an ideal of
    $\fsl(V_1)$, because $\fh$ surjects onto $\fsl(V_1)$. As $\fsl_2$ is
    simple, either it is all of $\fsl(V_1)$, in which case
    $\fh = \fsl(V_1) \oplus \fsl(V_2)$, or it is $0$. Arguing the same way
    with the second factor, if $\fh \neq \fsl(V_1) \oplus \fsl(V_2)$ then
    $\fh$ is the graph of an isomorphism $\fsl(V_1) \to \fsl(V_2)$ defined
    over $\Q$. Every such isomorphism has the form $x \mapsto f x f^{-1}$ for
    a linear isomorphism $f \colon V_1 \to V_2$, since
    $\Aut(\fsl_2) = \mathrm{PGL}_2$ and
    $\mathrm{PGL}_2(\Q) = \GL_2(\Q)/\Q^\times$. Then
    $\Hg(A) = \{(g, f g f^{-1})\}$ fixes $f \in \Hom(V_1, V_2)$, so $f$ is a
    morphism of Hodge structures by \Cref{cor:key}(c). This gives an
    isogeny $E_2 \to E_1$, a contradiction.

    For general $n$, induct: by induction $\fh$ surjects onto
    $\fsl(V_2) \oplus \cdots \oplus \fsl(V_n)$, and if
    $\fh \cap (\fsl(V_1) \oplus 0) = 0$, then $\fh$ is the graph of a surjection
    $\bigoplus_{i \geq 2} \fsl(V_i) \to \fsl(V_1)$. Its kernel is an ideal,
    hence a sum of all but one factor $\fsl(V_j)$, so the projection of $\fh$
    to $\fsl(V_1) \oplus \fsl(V_j)$ is a graph, contradicting the case
    $n = 2$.
\end{proof}

Note that $\SL_2 = \Sp_2$, so $\MT(E_1 \times \cdots \times E_n) =
\Gm \cdot (\Sp_{2,\Q})^n$. This is the shape that will appear in
\Cref{thm:main}.

\section{Hodge loci}\label{sec:hodge-loci}

\subsection{Definitions}
Let $f \colon X \to S$ be a smooth projective morphism to a connected complex
manifold, and fix $p$. Let $\pi \colon \wt{S} \to S$ be the universal cover.
Over $\wt{S}$ the local system $R^{2p}f_*\Q$ is trivial, so parallel transport
identifies $H^{2p}(X_{\tilde{s}}, \Q)$ with a fixed $\Q$-vector space $H$ for
all $\tilde{s} \in \wt{S}$. The Hodge filtration $F^\bullet_{\tilde{s}}$ on
$H_\C$ varies holomorphically with $\tilde{s}$ (Griffiths).

\begin{definition}
    For $\lambda \in H$, the \emph{Hodge locus} (or \emph{Noether--Lefschetz
    locus}) of $\lambda$ is
    \begin{equation*}
        \NL(\lambda) = \{\tilde{s} \in \wt{S} : \lambda \in
        \Hdg^p(X_{\tilde{s}})\}.
    \end{equation*}
    Its image $\pi(\NL(\lambda)) \subseteq S$ is the set of $s$ such that some
    parallel transport of $\lambda$ to $X_s$ is a Hodge class. The \emph{Hodge
    locus of $f$} is
    \begin{equation*}
        \HL(f) = \bigcup_{\lambda \in H,\ \NL(\lambda) \neq \wt{S}}
        \pi(\NL(\lambda)) \subseteq S.
    \end{equation*}
    The same definitions make sense for any polarizable variation of Hodge
    structure, e.g.\ for tensor constructions of $R^k f_* \Q$.
\end{definition}

\begin{example}[Noether--Lefschetz]
    Let $S$ parametrize smooth surfaces $X_s \subseteq \PP^3$ of degree
    $d \geq 4$, and $p = 1$. By Lefschetz $(1,1)$, $\HL(f)$ is the locus of
    surfaces with Picard number $\rho(X_s) \geq 2$. The Noether--Lefschetz
    theorem says that $\HL(f) \neq S$, i.e.\ the very general such surface has
    $\Pic(X_s) = \Z \cdot \cO(1)$. However, $\HL(f)$ is dense in $S$ in the
    analytic topology \cite{CHM88}. This is where the name comes from; for
    classes of higher degree, ``Hodge locus'' is more common.
\end{example}

\begin{proposition}
    $\NL(\lambda)$ is a closed complex analytic subset of $\wt{S}$.
\end{proposition}
\begin{proof}
    Since $\lambda$ is real and $F^p \cap \ol{F^p} = H^{p,p}$ in weight $2p$,
    $\lambda$ is a Hodge class at $\tilde{s}$ if and only if
    $\lambda \in F^p_{\tilde{s}}$. So $\NL(\lambda)$ is the zero locus of the
    image of the constant (hence holomorphic) section $\lambda$ in the
    holomorphic vector bundle $(H \otimes \cO_{\wt{S}})/F^p$.
\end{proof}

Much more is true.

\begin{theorem}[Cattani--Deligne--Kaplan \cite{CDK95}]
    If $S$ is quasi-projective and $f$ is algebraic, then $\pi(\NL(\lambda))$
    is a closed algebraic subvariety of $S$ for every $\lambda$.
\end{theorem}

\begin{remark}[Hodge conjecture $\Rightarrow$ CDK]
    Suppose the Hodge conjecture holds. Then $\lambda_s$ is Hodge if and only
    if it is a $\Q$-linear combination of classes of subvarieties of $X_s$. The
    relative Hilbert scheme of $X/S$ has countably many components, each
    projective over $S$, and the cohomology class of a subvariety is locally
    constant in flat families. So the locus where $\lambda$ is algebraic is a
    countable union of images of (products of) such components, hence a
    countable union of closed algebraic subvarieties of $S$. Thus CDK is
    exactly what the Hodge conjecture predicts, and is one of the best pieces
    of evidence for it.
\end{remark}

\begin{definition}
    A point $s \in S$ is \emph{very general} if $s \notin \HL(f)$.
\end{definition}

If $\NL(\lambda) \neq \wt{S}$, then $\NL(\lambda)$ is nowhere dense in
$\wt{S}$, so $\pi(\NL(\lambda))$ is a countable union of nowhere dense subsets
of $S$, and by CDK (when $S$ is quasi-projective) it is a proper closed
subvariety. There are countably many
$\lambda \in H$, so the very general points are the complement of a countable
union of proper closed subvarieties; by the Baire category theorem they are
dense. By construction:

\begin{lemma}
    If $s \in S$ is very general, then a class in $H^{2p}(X_s, \Q)$ is Hodge if
    and only if it remains Hodge under parallel transport to every fiber. In
    words: on a very general fiber, the only Hodge classes are the ones forced
    by the family.
\end{lemma}

\subsection{Hodge loci in $\cA_g$}
Over $\fH_g$ everything is explicit, since $H^{2p}(A_h, \Q) = \Lambda^{2p} V$
for every $h$. For a tensor $t \in T^{a,b}(V)$ (e.g.\ $t \in \Lambda^{2p}V$),
put
\begin{equation*}
    \NL(t) = \{h \in \fH_g : t \text{ is a Hodge class for } h\}.
\end{equation*}
For a $\Q$-algebraic subgroup $H \subseteq \Sp(V)$, put
\begin{equation*}
    \fH_g(H) = \{h \in \fH_g : h(U(1)) \subseteq H(\R)\} =
    \{h \in \fH_g : \Hg(A_h) \subseteq H\}.
\end{equation*}

\begin{proposition}
    \begin{enumerate}[label=(\roman*)]
        \item $\NL(t) = \fH_g(G_t)$, where $G_t = \Stab_{\Sp(V)}(t)$.
        \item $\fH_g(H)$ is stable under conjugation by $H(\R)$. In particular,
        it contains the orbit $H(\R) \cdot h$ of any of its points.
    \end{enumerate}
\end{proposition}
\begin{proof}
    (i) is \Cref{prop:key}, together with the fact that
    $h(U(1)) \subseteq \Sp(V_\R)$ for all $h \in \fH_g$. For (ii), if
    $h(U(1)) \subseteq H(\R)$ and $g \in H(\R)$, then
    $g h(U(1)) g^{-1} \subseteq H(\R)$.
\end{proof}

\begin{fact}[see \cite{Moonen04, Milne05}]\label{fact:special}
    Let $h \in \fH_g$ and $H = \Hg(A_h)$. Then $H(\R)^+ \cdot h$ is a closed
    complex submanifold of $\fH_g$ (a Hermitian symmetric domain), and its
    image in $\cA_g$ is a closed algebraic subvariety, called a \emph{special}
    (or \emph{Shimura}) subvariety. Its very general point has Hodge group
    exactly $H$. The irreducible components of Hodge loci in $\cA_g$ are
    exactly the special subvarieties, and there are countably many of them
    since there are countably many $\Q$-subgroups of $\Sp(V)$.
\end{fact}

So we have a dictionary:
\begin{center}
\renewcommand{\arraystretch}{1.3}
\begin{tabular}{c|c}
    Hodge theory & Group theory \\
    \hline
    Hodge class $t$ on $A_h$ & $\Q$-subgroup $G_t \subseteq \Sp(V)$ fixing $t$\\
    Hodge locus $\NL(t)$ & union of orbits $H(\R)^+ \cdot h$, $H \subseteq G_t$\\
    very general point of a special subvariety & $\Hg(A_h) = H$\\
    very general point of $\cA_g$ & $\Hg(A_h) = \Sp(V)$ (\Cref{thm:very-general})
\end{tabular}
\end{center}

\begin{example}[$\cA_1$]\label{ex:A1}
    By \Cref{prop:elliptic}, the possible Hodge groups are $\SL_2$ and the
    tori $U_K$. Since $U_K(\R) \cong U(1)$ is compact, the orbit
    $U_K(\R) \cdot h$ is the single point $\{h\}$. So the special subvarieties
    of $\cA_1$ are $\cA_1$ itself and the CM points, and the Hodge loci
    (e.g.\ of the classes on $E \times E$ from \Cref{ex:ExE}) are the CM
    points. These are countable and dense. So in theorems about the Hodge
    classes on a ``generic'' member of a family, \emph{general} (outside a
    proper Zariski closed subset) cannot be improved to anything better than
    \emph{very general}.
\end{example}

\begin{example}[$\cA_2$]\label{ex:A2}
    By Lefschetz $(1,1)$, the Hodge loci of classes in $H^2(A, \Q)$ are the
    loci where $\rho(A) \geq 2$; their components are called \emph{Humbert
    surfaces}. There are two kinds of generic Hodge group.
    \begin{enumerate}[label=(\alph*)]
        \item Split $V = V_1 \oplus V_2$ into $\psi$-orthogonal planes, so
        $\theta = \theta_1 + \theta_2$ with $\theta_i \in \Lambda^2 V_i$. Then
        $G_{\theta_1} = \SL(V_1) \times \SL(V_2)$, and
        $\NL(\theta_1) = \fH_1 \times \fH_1$, whose image in $\cA_2$
        is the locus of products $E_1 \times E_2$. Its generic Hodge group is
        $\SL_2 \times \SL_2 = (\Sp_{2,\Q})^2$ by
        \Cref{prop:product-elliptic}.
        \item Let $F$ be a real quadratic field acting on $V$, making $V$ a
        $2$-dimensional $F$-vector space with $\psi = \Tr_{F/\Q} \circ \psi_F$.
        Then $H = \Res_{F/\Q} \SL_2$ has $H(\R) = \SL_2(\R) \times \SL_2(\R)$,
        the orbit is again $\fH_1 \times \fH_1$, and its image in $\cA_2$ is a
        Hilbert modular surface, parametrizing abelian surfaces with real
        multiplication by $F$.
    \end{enumerate}
    In both cases $H_\C \cong \SL_{2,\C} \times \SL_{2,\C}$ acting on $V_\C$ by
    the sum of the two standard representations, but only in case (a) is $H$ a
    product over $\Q$.
\end{example}

\section{The very general abelian variety}

\begin{lemma}\label{lem:sp-invariants}
    Let $W$ be a $\Q$-vector space of dimension $2h$ and
    $\theta \in \Lambda^2 W$ nondegenerate. Then
    \begin{equation*}
        (\Lambda^\bullet W)^{\Sp(W)} = \Q[\theta]/(\theta^{h+1}).
    \end{equation*}
\end{lemma}
\begin{proof}
    By the Lefschetz decomposition (which is pure linear algebra here),
    $\Lambda^k W = \bigoplus_j \theta^j \wedge P^{k-2j}$, where
    $P^i \subseteq \Lambda^i W$ ($0 \leq i \leq h$) is the primitive subspace.
    This decomposition is $\Sp(W)$-equivariant since $\theta$ is invariant. For
    $1 \leq i \leq h$, $P^i_\C$ is a nontrivial irreducible representation of
    $\Sp_{2h}(\C)$ \cite[Theorem 17.5]{FH91}, so it has no invariants. Only
    $P^0 = \Q$ contributes, and its copy in $\Lambda^{2j} W$ is spanned by
    $\theta^j$.
\end{proof}

\begin{theorem}\label{thm:very-general}
    Let $(A, \theta)$ be a very general principally polarized abelian variety
    of dimension $g$. Then
    \begin{equation*}
        \Hg(A) = \Sp(V) \cong \Sp_{2g}, \qquad \MT(A) = \GSp(V), \qquad
        \Hdg^p(A) = \Q \cdot \theta^p.
    \end{equation*}
    In particular, $A$ satisfies the Hodge conjecture.
\end{theorem}
\begin{proof}
    For each of the countably many tensors $t \in T^{a,b}(V)$, $\NL(t)$ is a
    closed analytic subset of $\fH_g$. Take $h$ outside every $\NL(t)$ that is
    not all of $\fH_g$. (Note that $\gamma \cdot \NL(t) = \NL(\gamma t)$ for
    $\gamma \in \Sp_{2g}(\Z)$, so this is a condition on the image of $h$ in
    $\cA_g$.)

    \emph{Claim: every Hodge tensor $t$ for $h$ is $\Sp(V)$-invariant.}
    Indeed, $\NL(t) = \fH_g$, so $t$ is fixed by $h'(U(1))$ for every
    $h' \in \fH_g$. Since $\fH_g = \{g h g^{-1} : g \in \Sp(V_\R)\}$, this
    means $\Stab(t)$ contains every $\Sp(V_\R)$-conjugate of $h(U(1))$. These
    generate a nontrivial connected normal subgroup of $\Sp(V_\R)$, whose
    Zariski closure is all of $\Sp(V_\R)$ because $\Sp_{2g}$ is simple. So
    $\Sp(V) \subseteq \Stab(t)$.

    Now by \Cref{prop:key}, \Cref{cor:key}(a), and the claim,
    \begin{equation*}
        (\Lambda^{2p} V)^{\Sp(V)} \subseteq (\Lambda^{2p} V)^{\Hg(A)} =
        \Hdg^p(A) \subseteq (\Lambda^{2p} V)^{\Sp(V)},
    \end{equation*}
    and the outer terms are $\Q \theta^p$ by \Cref{lem:sp-invariants}. The class
    $\theta^p$ is algebraic: it is the class of
    $\Theta_1 \cap \cdots \cap \Theta_p$ for general translates $\Theta_i$ of a
    theta divisor. Finally, by \Cref{fact:deligne}, $\Hg(A)$ is the subgroup
    fixing all Hodge tensors, which by the claim contains $\Sp(V)$; combined
    with \Cref{cor:key}(a), $\Hg(A) = \Sp(V)$.
\end{proof}

\begin{remark}
    By \Cref{ex:A1}, and similarly because CM points are dense in $\cA_g$, the
    words ``very general'' cannot be replaced by ``general''.
\end{remark}

\subsection{Back to integral classes}

\begin{corollary}
    For a very general principally polarized abelian variety $(A, \theta)$ of
    dimension $g$,
    \begin{equation*}
        \Hdg^p(A, \Z) = \Z \cdot \frac{\theta^p}{p!}.
    \end{equation*}
\end{corollary}
\begin{proof}
    In a symplectic basis,
    $\frac{\theta^p}{p!} = \sum_{|I| = p} \bigwedge_{i \in I}(e_i \wedge f_i)$,
    which lies in $\Lambda^{2p} H^1(A, \Z)$ and is indivisible. So
    $\Q\theta^p \cap H^{2p}(A, \Z) = \Z \cdot \theta^p/p!$.
\end{proof}

So the integral Hodge conjecture for $A$ in degree $p$ is equivalent to
$\theta^p/p!$ being algebraic (we already know $\theta^p = p! \cdot
\theta^p/p!$ is). For $p = 1$ this is clear. For curves, $p = g - 1$, the class
$\theta^{g-1}/(g-1)!$ is called the \emph{minimal class}.
\begin{itemize}
    \item If $(A, \theta) = (J(C), \theta)$ is a Jacobian, then the
    Abel--Jacobi image of $C$ has class $\theta^{g-1}/(g-1)!$ (Poincar\'e's
    formula). Every principally polarized abelian variety of dimension
    $g \leq 3$ is a product of Jacobians, so the minimal class is algebraic
    there.
    \item Conversely, by the Matsusaka--Ran criterion
    \cite{Matsusaka59, Ran81}, if the minimal class is represented by an
    \emph{effective} curve then $(A, \theta)$ is a product of Jacobians. Since
    $\dim \mathcal{M}_g = 3g - 3 < g(g+1)/2 = \dim \cA_g$ for $g \geq 4$, the
    minimal class on a very general $A$ of dimension $\geq 4$ is not effective.
    It could still be a difference of effective classes.
    \item For $g = 4, 5$ the very general principally polarized abelian variety
    is a Prym variety, and the Abel--Prym curve has class
    $2\theta^{g-1}/(g-1)!$.
\end{itemize}

\begin{theorem}[Engel--de Gaay Fortman--Schreieder \cite{EdGFS25}]\label{thm:EdGFS}
    Let $(A, \theta)$ be a very general principally polarized abelian variety
    of dimension $g \geq 4$. Then the class of every curve on $A$ is an
    \emph{even} multiple of $\theta^{g-1}/(g-1)!$. In particular, the minimal
    class is a non-algebraic integral Hodge class, and the integral Hodge
    conjecture fails for $A$.
\end{theorem}

They prove the same for the intermediate Jacobian of a very general cubic
threefold, which implies that very general cubic threefolds are not stably
rational.

\begin{remark}[Exceptional Hodge classes]
    The very general abelian variety has as few Hodge classes as possible.
    Special abelian varieties can have Hodge classes that are not polynomials
    in divisor classes, and for these the Hodge conjecture is a real problem.
    The classical example is due to Weil \cite{Weil77}: if an imaginary
    quadratic field $K$ acts on an abelian fourfold $A$ so that each embedding
    of $K$ occurs with multiplicity $2$ on $H^{1,0}$, then
    $\Lambda^4_K H^1(A, \Q) \subseteq H^4(A, \Q)$ is a $2$-dimensional space of
    Hodge classes (\emph{Weil classes}), which for very general such $A$ are
    not generated by divisors. Proving that Weil classes are algebraic is a
    much-studied test case of the Hodge conjecture. See \cite{Gordon99} for a
    survey.
\end{remark}

\section{Lagrangian fibrations}

\begin{definition}[\cite{BakkerLehn22}]
    A \emph{primitive symplectic variety} is a normal compact K\"ahler space
    $X$ with a holomorphic symplectic form $\sigma$ on $X^{\mathrm{reg}}$ that
    extends holomorphically to a resolution of $X$, such that
    $H^1(X, \cO_X) = 0$ and $H^0(X, \Omega_X^{[2]}) = \C \sigma$. The smooth
    ones are exactly the irreducible hyperk\"ahler manifolds.
\end{definition}

Let $X$ be a primitive symplectic variety of dimension $2n$. A
\emph{Lagrangian fibration} is a surjective morphism $f \colon X \to B$ with
connected fibers onto a normal projective variety with
$0 < \dim B < 2n$. Then $\dim B = n$, the fibers are Lagrangian, and the fiber
$X_b$ over a general point $b \in B$ is an abelian variety of dimension $n$
(Matsushita \cite{Matsushita99}, in the smooth case).

Let $B^0 \subseteq B$ be the dense open subset over which $f$ is smooth. Fixing
an ample class on $X$ and restricting it to the fibers polarizes every $X_b$,
$b \in B^0$, with the same polarization type $\delta$. So we get a moduli map
\begin{equation*}
    \phi \colon B^0 \to \cA_{n, \delta}, \qquad b \mapsto [X_b],
\end{equation*}
where $\cA_{n,\delta}$ is the moduli space of $n$-dimensional abelian
varieties with a polarization of type $\delta$ (the analogue of $\cA_n$ for
non-principal polarizations).

\begin{definition}
    The Lagrangian fibration $f$ is \emph{maximally varying} if $\phi$ is
    generically finite onto its image, i.e.\ the isomorphism class of $X_b$
    varies in all $n = \dim B$ directions.
\end{definition}

\begin{example}\label{ex:K3n}
    Let $S \to \PP^1$ be an elliptic K3 surface whose $j$-invariant is not
    constant. Then
    \begin{equation*}
        S^{[n]} \to \Sym^n S \to \Sym^n \PP^1 \cong \PP^n
    \end{equation*}
    is a Lagrangian fibration, whose fiber over $b = \{t_1, \ldots, t_n\}$
    (distinct points with smooth fibers $E_{t_i}$) is
    $E_{t_1} \times \cdots \times E_{t_n}$. Since $j$ is not constant, this
    fibration is maximally varying. Since moreover isogeny classes are
    countable, for very general $b$ the curves $E_{t_i}$ have no CM and are
    pairwise non-isogenous. So by
    \Cref{prop:product-elliptic},
    \begin{equation*}
        \MT(X_b) = \Gm \cdot (\Sp_{2,\Q})^n.
    \end{equation*}
\end{example}

The main theorem says that this is what always happens, up to replacing
$\Sp_2$ by $\Sp_{2h}$.

\begin{theorem}[Varvak]\label{thm:main}
    Let $X$ be a primitive symplectic variety with a maximally varying
    Lagrangian fibration $f \colon X \to B$. For very general $b \in B$,
    \begin{equation*}
        \MT(X_b) \cong \Gm \cdot (\Sp_{2h, \Q})^k
    \end{equation*}
    for some $h, k \geq 1$. Equivalently, $\Hg(X_b) \cong (\Sp_{2h,\Q})^k$.
\end{theorem}

(For the equivalence: if $\MT = \Gm \cdot S$ with $S$ semisimple, then
$S = [\MT, \MT] = [\Hg, \Hg] \subseteq \Hg$, and $\Hg$ is a connected subgroup of
$(\MT \cap \SL(V))$, whose identity component is $S$.)

\begin{remark}[The maximal variation hypothesis is necessary]
    Let $A = E \times F$ be a product of elliptic curves and
    $X = K_n(A) \subseteq A^{[n+1]}$ the generalized Kummer variety. The
    projection $A \to E$ induces a Lagrangian fibration $X \to \PP^n$, and
    over a dense open subset of $\PP^n$ every fiber is isomorphic to
    \begin{equation*}
        \Big\{(y_1, \ldots, y_{n+1}) \in F^{n+1} : \sum_i y_i = 0\Big\}
        \cong F^n.
    \end{equation*}
    So this fibration is isotrivial ($\phi$ is constant). If $F$ has CM by $K$,
    then $\MT(X_b) = \MT(F) = \Res_{K/\Q}\Gm$ is a torus, not of the form in
    \Cref{thm:main}. (The Hodge conjecture still holds for $F^n$: the Hodge
    ring of a power of an elliptic curve is generated by divisor classes.)
    Matsushita conjectured that for a Lagrangian fibration of a hyperk\"ahler
    manifold, $\phi$ is either constant or generically finite onto its image,
    so conjecturally the isotrivial fibrations are the only ones excluded.
\end{remark}

\begin{corollary}\label{cor:main}
    In the situation of \Cref{thm:main}, for very general $b \in B$ the
    abelian variety $X_b$, and every power $X_b^m$, satisfies the Hodge
    conjecture.
\end{corollary}

The proof of \Cref{thm:main} is not the topic of this talk. Instead, let us
see how the corollary follows. It is a statement purely about abelian
varieties:

\begin{theorem}\label{thm:HC-from-MT}
    Let $A$ be an abelian variety with $\Hg(A) \cong (\Sp_{2h, \Q})^k$. Then
    every Hodge class on $A$ is a polynomial in divisor classes. In
    particular, $A$ satisfies the Hodge conjecture, and so does every power
    $A^m$.
\end{theorem}

The proof has two steps. First we pin down how $\Hg(A)$ acts on
$V = H^1(A, \Q)$, and then we compute invariants.

\begin{lemma}\label{lem:rep}
    Write $\Hg(A) = S_1 \times \cdots \times S_k$ with $S_i \cong \Sp_{2h,\Q}$,
    and let $W_i$ be the standard $2h$-dimensional representation of $S_i$.
    Let $V^{(i)} = V^{\prod_{j \neq i} S_j}$. Then
    \begin{equation*}
        V = \bigoplus_{i=1}^k V^{(i)}, \qquad V^{(i)} \cong W_i \otimes M_i,
    \end{equation*}
    where $M_i$ is a vector space with trivial action, of dimension $m_i$.
\end{lemma}
\begin{proof}
    Let $\mu = h|_{U(1)} \colon U(1) \to \Hg(A)(\R)$, with components
    $\mu_i \colon U(1) \to S_i(\R)$.
    \begin{enumerate}[label=(\alph*)]
        \item \emph{Each $\mu_i$ is nontrivial.} Otherwise
        $h(U(1)) \subseteq \prod_{j \neq i} S_j$, a smaller $\Q$-subgroup,
        contradicting the minimality of $\Hg(A)$. (This is where we use that
        the factors $S_i$ are defined over $\Q$.)
        \item \emph{Each irreducible constituent of $V_\C$ is a representation
        of a single factor.} Let $U = U_1 \boxtimes \cdots \boxtimes U_k$ be an
        irreducible constituent. The torus $\mu(U(1))$ acts on $V_\C$ with
        weights exactly $\pm 1$ (by $z$ on $V^{1,0}$ and $\ol{z} = z^{-1}$ on
        $V^{0,1}$). If $U_i$ is nontrivial, then $\mu_i$ acts on $U_i$
        nontrivially (as $S_i \to \GL(U_i)$ has finite kernel) with
        determinant $1$, so with weights $a < 0 < b$ among others. If $U_i$
        and $U_j$ were both nontrivial for some $i \neq j$, with $\mu_j$-weights
        $c < 0 < d$, then $\mu$ would have the three distinct weights
        $a + c < a + d < b + d$ on $U$. This is impossible, and $U$ cannot be
        trivial (weight $0$), so $U = U_i$ for exactly one $i$, and $\mu_i$
        acts on $U_i$ with weights exactly $\pm 1$.
        \item \emph{$U_i$ is the standard representation.} Since $\mu_i$ has
        weights $\pm 1$ on $U_i$, it has weights in $\{0, \pm 2\}$ on
        $\fsp_{2h} \subseteq \End(U_i)$. In standard coordinates, with $\mu_i$
        dominant, this forces $\mu_i = (1, \ldots, 1)$. If $\lambda$ is the
        highest weight of $U_i$, then the weights $\lambda$ and
        $(-\lambda_1, \lambda_2, \ldots, \lambda_h)$ pair with $\mu_i$ to
        $\sum_j \lambda_j \in \{\pm 1\}$ and
        $\sum_j \lambda_j - 2\lambda_1 \in \{\pm 1\}$, which forces
        $\lambda = (1, 0, \ldots, 0)$.
        \item So $V^{(i)}_\C$ is isotypic of type $W_{i,\C}$ and
        $V = \bigoplus_i V^{(i)}$. Since $W_i$ is defined over $\Q$ and
        absolutely irreducible,
        $V^{(i)} \cong W_i \otimes \Hom_{S_i}(W_i, V^{(i)})$. \qedhere
    \end{enumerate}
\end{proof}

\begin{remark}
    By \Cref{cor:key}, the $V^{(i)}$ are sub-Hodge structures and
    $\End_{\mathrm{HS}}(V^{(i)}) = \End(M_i)$. So
    $A \sim A_1^{m_1} \times \cdots \times A_k^{m_k}$ up to isogeny, with the
    $A_i$ pairwise non-isogenous abelian varieties of dimension $h$ satisfying
    $\Hg(A_i) = \Sp_{2h}$, and $\dim A = h \sum_i m_i$. In
    \Cref{ex:K3n}, $h = 1$, $k = n$, and all $m_i = 1$.
\end{remark}

\begin{remark}[Why splitting over $\Q$ matters]
    The abstract group $\Hg(A)_\C$ does not determine the Hodge classes.
    Mumford \cite{Mumford69} constructed abelian fourfolds with
    $\End(A) = \Z$ and $\Hg(A)_\C \cong \SL_{2,\C}^3$, acting on
    $V_\C \cong \C^2 \otimes \C^2 \otimes \C^2$. There $\Hg(A)$ is a non-split
    $\Q$-form of $\SL_2^3$, two of the three factors of $\Hg(A)_\R$ are
    compact, $\mu$ is trivial on them, and step (a) above fails. Similarly,
    \Cref{ex:A2}(b) shows that $\Hg(A)_\C \cong \Sp_{2h,\C}^k$ does not imply
    $\Hg(A) \cong \Sp_{2h,\Q}^k$.
\end{remark}

The second step is classical invariant theory.

\begin{lemma}[First fundamental theorem for $\Sp$, skew version \cite{Howe89}]\label{lem:FFT}
    Let $W$ be a symplectic $\Q$-vector space with symplectic class $\theta$,
    and $M$ a finite-dimensional $\Q$-vector space with trivial
    $\Sp(W)$-action. Then the algebra $(\Lambda^\bullet(W \otimes M))^{\Sp(W)}$
    is generated by its degree $2$ part, which is
    $\theta \otimes \Sym^2 M \subseteq (\Lambda^2 W \otimes \Sym^2 M) \oplus
    (\Sym^2 W \otimes \Lambda^2 M) = \Lambda^2(W \otimes M)$.
\end{lemma}

When $\dim M = 1$ this is just \Cref{lem:sp-invariants}.

\begin{proof}[Proof of \Cref{thm:HC-from-MT}]
    By \Cref{lem:rep},
    \begin{equation*}
        H^\bullet(A, \Q) = \Lambda^\bullet V =
        \bigotimes_{i=1}^k \Lambda^\bullet (W_i \otimes M_i),
    \end{equation*}
    and $S_i$ acts only on the $i$-th factor. Hence by \Cref{prop:key},
    \begin{equation*}
        \Hdg^\bullet(A) = \big(\Lambda^\bullet V\big)^{S_1 \times \cdots \times
        S_k} = \bigotimes_{i=1}^k \big(\Lambda^\bullet(W_i \otimes
        M_i)\big)^{S_i}.
    \end{equation*}
    By \Cref{lem:FFT}, this algebra is generated by its degree $2$ part, which
    consists of Hodge classes in $H^2(A, \Q)$. By Lefschetz $(1,1)$ these are
    classes of $\Q$-divisors, and products of divisor classes are algebraic
    (intersect general translates). So every Hodge class on $A$ is algebraic.

    For $A^m$: by \Cref{rmk:MT-basic}(4), $\Hg(A^m) = \Hg(A)$, now acting on
    $V^{\oplus m} = \bigoplus_i W_i \otimes (M_i \otimes \Q^m)$, and the same
    argument applies.
\end{proof}

\begin{remark}
    When all $m_i = 1$ (as in \Cref{ex:K3n}), only \Cref{lem:sp-invariants} is
    needed: the Hodge ring of $A$ is $\bigotimes_i \Q[\theta_i]/(\theta_i^{h+1})$,
    where $\theta_i \in \Lambda^2 W_i \subseteq H^2(A, \Q)$ is the pullback of a
    polarization along $A \to A_i$. This is exactly parallel to
    \Cref{thm:very-general}, with $\Sp_{2h}^k$ in place of $\Sp_{2g}$.
\end{remark}


\bibliographystyle{alpha}
\bibliography{refs}


\end{document}
